\documentclass[10pt,a4paper]{amsart}
\usepackage[margin=19mm]{geometry}
\usepackage{amsmath,amssymb,mathtools}
\usepackage[scr=rsfs,cal=euler]{mathalfa}
\usepackage{xfrac}
\usepackage[unicode,hidelinks,bookmarks=false]{hyperref}
\newcommand{\ie}{i.e.\ }
\newcommand{\cf}{cf.\ }
\newcommand{\resp}{respectively\ }
\newcommand{\Subsection}{\subsection}
\providecommand{\nobreakdash}{\nobreak\hbox{-}}
\setlength{\emergencystretch}{2em}
\multlinegap0pt
\usepackage{bm}
\usepackage{bbm}
\usepackage{dsfont}
\usepackage{xspace}
\usepackage{cancel}
\usepackage{mathtools}
\usepackage{amsthm}
\makeatletter
\@ifundefined{theorem}{\newtheorem{theorem}{Theorem}}{}
\@ifundefined{defn}{\newtheorem{defn}[theorem]{Definition}}{}
\@ifundefined{lemma}{\newtheorem{lemma}[theorem]{Lemma}}{}
\makeatother
\title[The Game of Graph Nim on graphs with four edges]{The Game of Graph Nim on graphs with four edges}
\author{Sayar Karmakar, Moumanti Podder, Souvik Roy, Soumyarup Sadhukhan}
\date{Source: arXiv:2509.05064v1 (CC BY 4.0)}
\begin{document}
\maketitle

\noindent\emph{Source and licence.} The Game of Graph Nim on graphs with four edges. Version arXiv:2509.05064v1 (CC BY 4.0), distributed under the Creative Commons Attribution 4.0 International licence, as recorded in the arXiv licence metadata for this exact version and shown on the abstract page https://arxiv.org/abs/2509.05064v1. Full rights notice: licenses/arxiv-2509.05064-CC-BY-4.0.txt.\par\smallskip

\noindent\textit{Editorial scope.} Counted unit: the Lemma whose source label is \texttt{lem:nec\_suff} in the article (source0.tex lines 579--584, source label \texttt{lem:nec\_suff}), delivered together with the article's own complete original proof of it (source0.tex lines 585--604, one \texttt{proof} environment of 20 lines). No mathematics is written, repaired or re-derived: statement and proof are the article's own text line for line. Required context retained uncounted: defn:special (lines 206--212); M\_G\_{4 (lines 208--210). Every definition, display and label of those lines is retained unchanged. Omitted from this package and not redistributed: 1--205, 211--578, 605--2559. Nothing inside a retained excerpt is omitted or altered: every retained body is delivered exactly as the article's lines are written, so no omission receipt of any kind accompanies this package. Citations: the retained excerpts carry no citation command at all, so no bibliography entry of the article is transcribed and no reference is redistributed. Reference adaptation: none was needed and none was made; every reference inside the delivered unit resolves inside it, so no unresolved reference or citation is produced. Printed slips kept literal: no printed word, symbol, hypothesis or number of the counted statement or of its proof was altered. Layout and class: the article's own class is \texttt{amsart} with packages including mathrsfs, url, mathtools, latexsym, epsfig, amssymb, amsmath, amsthm, color, bm, enumitem, hyperref, amsaddr, amsbsy; the wrapper is the lane's pinned \texttt{amsart} editorial wrapper at 10pt on A4, which restates the theorem-like environments the retained text needs and adds the article's own macro definitions that the retained text needs (none), with extra wrapper packages (bm, bbm, dsfont, xspace, cancel, mathtools). The delivered numbering is the unit's own. Assets: no retained span contains a figure environment, a graphics-inclusion command, a PDF-page inclusion, a table or a TikZ picture; no image file is copied into this package, the package declares no asset dependency and no image file is redistributed with it. Licence: version arXiv:2509.05064v1 of this article is distributed under the Creative Commons Attribution 4.0 International licence, as recorded in the arXiv licence metadata for this exact version and shown on the abstract page https://arxiv.org/abs/2509.05064v1. A full rights notice is delivered at licenses/arxiv-2509.05064-CC-BY-4.0.txt. Characters: every retained excerpt is pure ASCII, so the delivered engine, XeLaTeX with its default Latin Modern text font, reports no missing character.
\begin{defn}\label{defn:special}
Given $k, m, i \in \mathbb{N}$ and $\ell \in \mathbb{N}_{0}$ such that
\begin{equation}\label{M_G_{4}}
i \in \{1,2,\ldots,m+1\} \quad \text{and} \quad \frac{m(m+1)}{2} \leqslant k \leqslant \frac{m(m+3)}{2},
\end{equation}
we say that a multiset $S$, consisting of $3$ elements in total, is \emph{$(k,\ell,m,i)$-special} if $S$, as a multiset, equals $\{k+1+(m+1)\ell,k+i+(m+1)\ell,k+m+2-i+(m+1)\ell\}$. We call a multiset $S$ simply \emph{special} if $S$ is $(k,\ell,m,i)$-special for \emph{some} $k, m, i \in \mathbb{N}$ and $\ell \in \mathbb{N}_{0}$ satisfying \eqref{M_G_{4}}.
\end{defn}

\begin{lemma}\label{lem:nec_suff}
Given $w_{0}(AB) \leqslant w_{0}(BC) \leqslant w_{0}(CA)$, with $w_{0}(AB)$, $w_{0}(BC)$ and $w_{0}(CA)$ not all equal, the multiset $\{w_{0}(AB),w_{0}(BC),w_{0}(CA)\}$ is special (see Definition~\ref{defn:special}) if and only if we have
\begin{equation}
2w_{0}(AB)>\left(w_{0}(BC)+w_{0}(CA)-2w_{0}(AB)\right)\left(w_{0}(BC)+w_{0}(CA)-2w_{0}(AB)+1\right).\label{eq:nec_suff}
\end{equation}
\end{lemma}

\begin{proof}
When \eqref{eq:nec_suff} holds, we set $m=w_{0}(BC)+w_{0}(CA)-2w_{0}(AB)$, and we let $k$ be the unique positive integer that, along with this choice of $m$, satisfies the second criterion of \eqref{M_G_{4}} as well as the relation $w_{0}(AB)-1\equiv k\bmod(m+1)$. Note that the inequality in \eqref{eq:nec_suff} ensures that 
\begin{equation}
w_{0}(AB)>\frac{m(m+1)}{2} \implies w_{0}(AB)\geqslant \frac{m(m+1)}{2}+1 \implies w_{0}(AB)-1\geqslant \frac{m(m+1)}{2},\nonumber
\end{equation}
so that, from the way we have defined $k$ above, we know that there exists $\ell \in \mathbb{N}_{0}$ such that $w_{0}(AB)=k+1+(m+1)\ell$. Next, from our definition of $m$, we have
\begin{align}
w_{0}(BC)+w_{0}(CA)=2w_{0}(AB)+m=2\{k+1+(m+1)\ell\}+m=2k+m+2+2(m+1)\ell.\label{nec_suff_eq_1}
\end{align}
Let $w_{0}(BC)-k\equiv i \bmod (m+1)$ for some $i \in \{1,2,\ldots,m+1\}$. Note that, since $w_{0}(BC)\geqslant w_{0}(AB)=k+1+(m+1)\ell$, we can find $\ell_{1} \in \mathbb{N}_{0}$ such that $w_{0}(BC)=k+i+(m+1)\ell_{1}$. Likewise, setting $w_{0}(CA)-k\equiv j \bmod(m+1)$ for some $j \in \{1,2,\ldots,m+1\}$, we see that there exists some $\ell_{2}\in\mathbb{N}_{0}$ such that $w_{0}(CA)=k+j+(m+1)\ell_{2}$. From \eqref{nec_suff_eq_1}, we then obtain
\begin{align}
{}&\{k+i+(m+1)\ell_{1}\}+\{k+j+(m+1)\ell_{2}\}=2k+m+2+2(m+1)\ell \nonumber\\
{}&\Longleftrightarrow i+j+(m+1)(\ell_{1}+\ell_{2})=m+2+2(m+1)\ell.\nonumber
\end{align}
The first conclusion to draw from this is that $(i+j) \equiv 1 \bmod (m+1)$, and since $2 \leqslant (i+j) \leqslant 2(m+1)$, the only possibility we are left with is that $(i+j)=(m+2)$. This, in turn, yields $(\ell_{1}+\ell_{2})=2\ell$. If possible, let us assume that $\ell_{1}\neq \ell_{2}$, which implies that $\ell_{1}<\ell_{2}$ (since $w_{0}(BC)\leqslant w_{0}(CA)$). This implies that $\ell_{1}\leqslant (\ell-1)$, which yields
\begin{align}
w_{0}(BC)=k+i+(m+1)\ell_{1}\leqslant k+i+(m+1)(\ell-1)=k+i-(m+1)+(m+1)\ell<k+1+(m+1)\ell=w_{0}(AB),\nonumber
\end{align}
leading to a contradiction. Therefore, our assumption that $\ell_{1}\neq\ell_{2}$ must be erroneous, leading to $\ell_{1}=\ell_{2}=\ell$. We have, therefore, deduced the existence of $k, m, i \in \mathbb{N}$ and $\ell \in \mathbb{N}_{0}$ such that \eqref{M_G_{4}} holds and $w_{0}(AB)=k+1+(m+1)\ell$, $w_{0}(BC)=k+i+(m+1)\ell$ and $w_{0}(CA)=k+j+(m+1)\ell=k+m+2-i+(m+1)\ell$. This proves one side of the implication stated in Lemma~\ref{lem:nec_suff}. The proof that the converse is also true is much more straightforward and hence omitted from this paper.
\end{proof}

\end{document}
